% Chapter 1: Tổng quan đề tài
\chapter{Tổng quan đề tài}

Máy bay bốn cánh (quadcopter) là một hệ thống khó điều khiển~\cite{quan2020}: phản ứng của máy bay phi tuyến (không tỷ lệ thuận với lệnh điều khiển), nhiều biến phải quản lý cùng lúc, và nó chỉ giữ được thăng bằng khi hai bài toán được giải liên tục với tốc độ cao: ước lượng tư thế (máy bay đang nghiêng bao nhiêu, quay về hướng nào) và điều khiển động cơ. Xuất phát từ đặc thù đó, đề tài ``Thuật toán điều khiển máy bay bốn cánh sử dụng PID và bộ lọc Kalman'' (mã SVOISP-2025-CK-57) chọn tự xây dựng bộ điều khiển bay (flight controller) \textbf{FC-v01} từ đầu thay vì dùng sẵn mã nguồn mở~\cite{ardupilot-esp32,px4-hw}, để nhóm chủ động toàn bộ chuỗi từ cảm biến, ước lượng tư thế cho tới điều khiển, trên một nền tảng đồng nhất. Nền tảng đó gồm: vi điều khiển ESP32-S3 -- vi mạch tính toán hai nhân hoạt động ở 240~MHz; hệ điều hành thời gian thực FreeRTOS tích hợp sẵn trong bộ công cụ ESP-IDF~\cite{idf-freertos}, bảo đảm các công việc quan trọng luôn được thực hiện đúng thời hạn; mô-đun cảm biến 10-DOF GY-87 (gia tốc kế, con quay hồi tiếp, từ kế, cảm biến áp suất); thẻ microSD ghi dữ liệu bay; và ứng dụng trạm mặt đất giám sát hệ thống không dây qua Bluetooth.

Đề tài tập trung vào ba mục tiêu. \textit{Thứ nhất}, phần cứng: nhóm thiết kế và chế tạo hai phiên bản bản mạch in (PCB) bốn lớp đồng, kích thước $50 \times 50$~mm, gắn trực tiếp lên khung máy bay phổ thông. Phiên bản v2 đang sử dụng tách hai đường nguồn 3,3~V -- một đường riêng cho khối cảm biến, một cho phần còn lại -- để nhiễu điện không làm sai lệch số đo của cảm biến. \textit{Thứ hai}, firmware, tức chương trình chạy trên vi điều khiển: phần mềm chia năm tầng, mỗi tầng chỉ được gọi tầng dưới nó; mười chín công việc chạy song song (tác vụ) do hệ điều hành phân công; dữ liệu trao đổi giữa các tác vụ qua các cấu trúc cấp phát một lần khi khởi động; và một máy trạng thái hữu hạn (bộ điều phối chỉ cho phép hệ thống chuyển giữa các trạng thái xác định) quản lý toàn bộ vòng đời. \textit{Thứ ba}, khối ước lượng trạng thái: bộ lọc Kalman mở rộng (EKF, thuật toán hợp nhất dữ liệu nhiều cảm biến, cân nhắc độ tin cậy từng nguồn đo) bảy trạng thái ước lượng tư thế ở 500~Hz (500 lần mỗi giây), một bộ lọc độ cao ba trạng thái, tăng tốc bằng lệnh SIMD (một lệnh xử lý nhiều số cùng lúc), với dữ liệu đầu vào đã qua hiệu chuẩn cảm biến.

Phương pháp nghiên cứu là phát triển lặp theo mô-đun: từng khối chức năng được làm và thử nghiệm độc lập trên phần cứng trước khi tích hợp. Trong phạm vi thời gian đề tài, chuỗi cảm biến -- tiền xử lý -- ước lượng trạng thái -- truyền và ghi dữ liệu đã hoàn thành và chạy ổn định ở mức mô-đun. Khối điều khiển -- dù là mục tiêu đăng ký -- chưa được hiện thực: qua đánh giá, PID bộc lộ nhiều hạn chế trên nền tảng này, nên nhóm chuyển hướng nghiên cứu INDI (Incremental Nonlinear Dynamic Inversion) cho giai đoạn tiếp theo; hướng đi này cùng kế hoạch tích hợp và bay thực nghiệm được trình bày ở chương kết luận.

Báo cáo gồm năm chương nội dung: Chương~2 trình bày thiết kế hệ điều hành thời gian thực (RTOS); Chương~3 thiết kế phần cứng bộ điều khiển bay; Chương~4 bộ lọc Kalman bảy trạng thái; Chương~5 kết quả thực nghiệm; Chương~6 kết luận và hướng phát triển.
